\documentclass[12pt,a4paper]{article}
\usepackage[utf8]{inputenc}
\usepackage{amsmath,amssymb,amsthm}
\usepackage{algorithm}
\usepackage{algorithmic}
\usepackage{geometry}
\geometry{margin=1in}

\newtheorem{theorem}{Theorem}
\newtheorem{lemma}[theorem]{Lemma}
\newtheorem{proposition}[theorem]{Proposition}
\theoremstyle{definition}
\newtheorem{definition}[theorem]{Definition}

\title{Problem 29: Distinct Powers}
\author{}
\date{}

\begin{document}
\maketitle

\section{Problem Statement}

How many distinct terms are in the sequence $\{a^b : 2 \leq a \leq 100, \; 2 \leq b \leq 100\}$?

\section{Mathematical Development}

\begin{definition}[Primitive Base]
An integer $c \geq 2$ is a \emph{primitive base} if there do not exist $d \geq 2$ and $k \geq 2$ with $c = d^k$.
\end{definition}

\begin{definition}[Canonical Representation]
For $a \geq 2$, the \emph{canonical representation} is the unique pair $(c, p)$ with $c$ primitive, $p \geq 1$, and $a = c^p$.
\end{definition}

\begin{theorem}[Collision Criterion]\label{thm:collision}
$a_1^{b_1} = a_2^{b_2}$ if and only if $a_1, a_2$ share the same primitive base~$c$ (say $a_i = c^{p_i}$) and $p_1 b_1 = p_2 b_2$.
\end{theorem}

\begin{proof}
$(\Leftarrow)$ Immediate: $c^{p_1 b_1} = c^{p_2 b_2}$.

$(\Rightarrow)$ Let $a_i = c_i^{p_i}$ canonically. Then $c_1^{p_1 b_1} = c_2^{p_2 b_2}$. By the fundamental theorem of arithmetic and primitivity of $c_1, c_2$, we obtain $c_1 = c_2$ and $p_1 b_1 = p_2 b_2$.
\end{proof}

\begin{theorem}[Reduction to Exponent Counting]\label{thm:reduction}
Let $K_c = \max\{p : c^p \leq 100\}$ for each primitive base~$c$, and define
\[
E_c = \{p \cdot b : 1 \leq p \leq K_c, \; 2 \leq b \leq 100\}.
\]
The number of distinct terms equals $\displaystyle\sum_{c \text{ primitive}} |E_c|$.
\end{theorem}

\begin{proof}
By Theorem~\ref{thm:collision}, terms from different base groups are always distinct, and terms within the same base group are distinct if and only if their effective exponents differ.
\end{proof}

\begin{lemma}[Base Groups with Collisions]\label{lem:groups}
The primitive bases with $K_c \geq 2$ are: $c \in \{2, 3, 5, 6, 7, 10\}$ with $K_c \in \{6, 4, 2, 2, 2, 2\}$ respectively. All other primitive bases have $K_c = 1$ and contribute exactly $99$ distinct terms each.
\end{lemma}

\begin{lemma}[Exponent Count for $K_c = 2$]\label{lem:k2}
For $K_c = 2$: $|E_c| = |[2,100] \cup [4,200]| = 99 + 99 - 49 = 149$.
\end{lemma}

\begin{proof}
$[2,100] \cap [4,200] = \{4, 6, 8, \ldots, 100\}$ has $49$ elements. Apply inclusion-exclusion.
\end{proof}

\begin{theorem}[Answer]\label{thm:answer}
The number of distinct terms is $9183$.
\end{theorem}

\begin{proof}
Sum $|E_c|$ over all primitive bases. The six non-trivial groups ($K_c \geq 2$) are computed by explicit exponent-set enumeration. All remaining bases contribute $99$ each. The total is verified by direct computation with arbitrary-precision arithmetic.
\end{proof}

\section{Algorithm}

\begin{algorithm}[H]
\caption{Distinct Powers via Canonical Representation}
\begin{algorithmic}[1]
\REQUIRE Upper bounds $A$, $B$
\ENSURE Number of distinct terms in $\{a^b : 2 \leq a \leq A, \; 2 \leq b \leq B\}$
\FOR{$a = 2$ \TO $A$}
    \STATE Compute canonical $(c, p)$ such that $a = c^p$
\ENDFOR
\STATE Initialize map $M : \text{base} \to \text{set of integers}$
\FOR{$a = 2$ \TO $A$}
    \STATE Let $(c, p) = \text{canonical}(a)$
    \FOR{$b = 2$ \TO $B$}
        \STATE Insert $p \cdot b$ into $M[c]$
    \ENDFOR
\ENDFOR
\RETURN $\sum_{c} |M[c]|$
\end{algorithmic}
\end{algorithm}

\section{Complexity Analysis}

\begin{proposition}
The canonical-representation method runs in $O(N^2)$ time and $O(N^2)$ space, where $N = A - 1$.
\end{proposition}

\begin{proof}
The double loop performs $N^2 = 9801$ set insertions. Each effective exponent is at most $K_c \cdot B \leq 600$, so all arithmetic is $O(1)$.
\end{proof}

\section{Answer}

\[
\boxed{9183}
\]

\end{document}
